% ============================================================================
% Appendix E.2: Positioning Among Bandit-Type LLM Routers
% ============================================================================
% Purpose: Detailed comparison with PILOT, BaRP, LLM Bandit, and simpler
%          bandit families. Supports Section 6 (Related Work) with expanded
%          taxonomy and key architectural differentiators.

\subsection{Positioning Among Bandit-Type LLM Routers}
\label{appendix:positioning}

The LLM routing landscape now includes multiple bandit-based approaches alongside
established supervised routers.
Table~\ref{tab:bandit_taxonomy} provides a structured comparison across five
families: simple MABs, plain contextual bandits, supervised static routers,
concurrent bandit routers, and banditGPT.

% --------------------------------------------------------------------------
% Taxonomy Table
% --------------------------------------------------------------------------

\begin{table*}[t]
\centering
\caption{Taxonomy of LLM routing approaches. \checkmark~= open-source code available.
$\dagger$~= code not publicly available.}
\label{tab:bandit_taxonomy}
\small
\begin{tabular}{@{}lp{3.0cm}p{2.8cm}p{3.8cm}p{4.0cm}@{}}
\toprule
\textbf{Aspect} & \textbf{Static Routers} & \textbf{Plain LinUCB} &
  \textbf{PILOT$^\dagger$ / BaRP$^\dagger$ / LLM Bandit} & \textbf{banditGPT (Ours)} \\
\midrule
Context use &
  Dense embedding (RouteLLM); graph (GraphRouter); contrastive (RouterDC) &
  Linear over prompt features &
  Shared query--LLM embeddings (PILOT); prompt + preference vector (BaRP); model ID vectors (LLM Bandit) &
  PCA-compressed prompt embeddings + bias via \texttt{FeatureService} \\
\midrule
Online adaptation &
  None (frozen at training time) &
  Yes (LinUCB updates) &
  Yes (all three adapt online from bandit feedback) &
  Yes (LinUCB updates + meta-learner reweighting) \\
\midrule
Cost awareness &
  Threshold (RouteLLM); edge weights (GraphRouter) &
  Ad-hoc penalties &
  Knapsack policy (PILOT); scalarized reward (BaRP); preference conditioning (LLM Bandit) &
  Hard constraints + cost-penalized LinUCB \\
\midrule
Prior machinery &
  Supervised offline training (100k+ labeled pairs) &
  Isotropic $\lambda I$, no warmup &
  Preference-informed from human data (PILOT); bandit-feedback trained (BaRP) &
  80k-battle warmup priors, heuristic priors, semantic neighbor transfer \\
\midrule
Safety / meta-learning &
  None &
  None &
  None &
  Corralling (Log-Barrier OMD) between warmup and tabula rasa experts \\
\midrule
Constraint handling &
  Implicit (threshold or classifier) &
  None or soft &
  Knapsack (PILOT); preference scalarization (BaRP) &
  Per-request hard \texttt{max\_cost}, \texttt{max\_latency},
  \texttt{quality\_floor} \\
\midrule
Multi-model ($K$) &
  $K{=}2$ (RouteLLM); $K{\geq}2$ (GraphRouter, RouterDC) &
  $K{\geq}2$ &
  $K{\geq}2$ (all three) &
  $K{\geq}2$ (evaluated at $K{=}2, 3$) \\
\midrule
Open source &
  RouteLLM \checkmark, GraphRouter \checkmark, RouterDC \checkmark &
  Varies &
  LLM Bandit \checkmark~(minimal); PILOT $\dagger$; BaRP $\dagger$ &
  \checkmark~(production-grade) \\
\bottomrule
\end{tabular}
\end{table*}

% --------------------------------------------------------------------------
% Key Differentiators
% --------------------------------------------------------------------------

\paragraph{1.\ Corralling Meta-Learner (Bandits over Bandits).}
PILOT, BaRP, LLM Bandit, and C2MAB-V~\cite{dai2024c2mab} each use a single
bandit policy or policy gradient to make routing decisions.
banditGPT adds a meta-learning layer: a Corralling~\cite{agarwal2017corralling}
coordinator that maintains trust weights over two expert bandits with
complementary exploration schedules.
The ``bandits over bandits'' structure provides a formal safety guarantee
(Appendix~\ref{appendix:mathematical_foundations}): even under severe prior
mismatch, total regret is bounded by tabula-rasa performance plus
$O(\sqrt{T})$ meta-overhead (Corollary~\ref{cor:neg_transfer}).
No concurrent bandit-based router provides an analogous meta-ensemble or
formal safety bound against prior mismatch.
Subsequent model-selection frameworks~\cite{pacchiano2020model, cutkosky2021dynamic}
improve on Corralling's variance guarantees through regret balancing, though
none have been applied to LLM routing.

\paragraph{2.\ Rich Prior and Registration Machinery.}
banditGPT maintains three prior layers: (a)~warmup priors from 80k RouteLLM
battles (or smaller domain-specific evaluation batches for $K{>}2$),
(b)~heuristic priors for models absent from the warmup data, and
(c)~semantic neighbor transfer with confidence reset for model onboarding.
New models are admitted via \texttt{register\_model()}, which copies $\theta$
from the nearest neighbor while resetting $A$ to $\neff \cdot I$, inheriting
preferences without inheriting over-confidence.
BaRP trains under bandit feedback; LLM Bandit uses model identity vectors.
Both are effective strategies but without the warm-start advantage that
pre-existing evaluation or preference data provides.

\paragraph{2a.\ Warmup Priors: banditGPT vs.\ PILOT.}
PILOT is the closest concurrent system to banditGPT in its use of offline
data to warm-start LinUCB, and the two approaches share the same Bayesian
intuition: encode offline knowledge into the sufficient statistics
$(A, b)$ so that the arm-parameter posterior begins near a useful estimate
rather than at the uninformative origin.
Despite this shared motivation, the mechanisms differ along four axes:
\begin{itemize}[nosep]
  \item \textbf{Prior signal.}
    banditGPT accumulates $(A_m, b_m)$ directly from (prompt, model, reward)
    triples via rank-1 LinUCB updates ($A_m \mathrel{+}= xx^\top$,
    $b_m \mathrel{+}= rx$), encoding absolute reward magnitudes in the
    context feature space.
    PILOT learns LLM embeddings $\theta_a^{\mathrm{pref}}$ from pairwise
    preference comparisons (which model ``won'') via triplet loss, then
    injects them as $b_a{=}\lambda_a\theta_a^{\mathrm{pref}}$ with
    $A_a{=}\lambda_a I$.
    Because banditGPT's priors are built from reward observations rather
    than relative comparisons, they directly encode the reward surface
    that the bandit will optimize, potentially providing a tighter initial
    estimate when reward magnitudes vary across models.
  \item \textbf{Prior strength control.}
    banditGPT uses a plasticity factor $\rho$ (default 0.1) that uniformly
    scales both $A$ and $b$, preserving the mean prediction while shrinking
    the effective sample size.
    PILOT sets $\lambda_a$ per arm as the inverse of the arm's pretraining
    accuracy, concentrating prior confidence on models whose offline
    performance is already well-estimated.
    banditGPT's uniform scaling is simpler but forgoes per-model calibration;
    PILOT's per-arm $\lambda_a$ is more expressive but requires access to
    per-arm accuracy metrics from the pretraining phase.
  \item \textbf{Embedding space.}
    PILOT jointly optimizes query and LLM embeddings in a shared space
    aligned by cosine affinity, potentially capturing non-linear
    query--model interactions.
    banditGPT uses a fixed SentenceTransformer + PCA pipeline with per-arm
    $(A_m, b_m)$ matrices, which is simpler to deploy and debug but cannot
    learn cross-model embedding structure.
  \item \textbf{Robustness to prior misspecification.}
    This is the most consequential difference.
    PILOT relies solely on $\lambda_a$ tuning and the natural overriding of
    the prior by online data; if the preference prior is misleading,
    recovery depends on the rate at which online observations dominate.
    banditGPT explicitly hedges via the Corralling meta-learner, which
    maintains a parallel tabula rasa expert and shifts weight to it within
    ${\sim}50$ decisions when warmup priors degrade
    (Section~\ref{sec:prior_degradation}).
    This architectural choice bounds negative transfer at
    $O(\sqrt{T \ln 2})$ meta-overhead
    (Corollary~\ref{cor:neg_transfer}), a guarantee absent from PILOT.
\end{itemize}

\paragraph{3.\ Production Engineering Around LinUCB.}
The \texttt{DisjointLinUCBPolicy} in \texttt{router.py} includes per-model
locks, a global lock, Sherman--Morrison $O(d^2)$ updates with automatic
$O(d^3)$ fallback when the denominator is small, forgetting factors,
staleness-aware variance inflation, and a numerical-stability pipeline
(trace-based checks, regularization resets, proactive regularization floor).
State save/load includes dimension validation and matrix-shape checks.
Thread-safe cloning uses \texttt{deepcopy} with atomic arm add/remove.
These engineering safeguards are essential for production deployments where
numerical instability or race conditions would degrade routing quality
silently.

\paragraph{4.\ Explicit Constraint Handling.}
\texttt{RouterConfig} carries \texttt{market\_cost\_floor/ceiling} and
\texttt{latency\_floor/ceiling}, plus pessimistic defaults when registry
metadata is missing (unknown models are treated as expensive and slow rather
than violating the budget).
Hard constraints (\texttt{max\_cost}, \texttt{max\_latency},
\texttt{quality\_floor}) are enforced via Dynamic Pareto Filtering
\emph{before} the bandit selects, so the bandit operates only over the
feasible set.
PILOT handles budget via a knapsack policy that selects a cost tier per
query; BaRP scalarizes cost into the reward; LLM Bandit uses preference
conditioning.
banditGPT's approach guarantees constraint satisfaction per-request rather
than in expectation over a horizon.

% --------------------------------------------------------------------------
% What Concurrent Work Does Better
% --------------------------------------------------------------------------

\paragraph{Advantages of Concurrent Work.}
We note areas where concurrent systems advance the state of the art beyond
banditGPT's current design:
\begin{itemize}[nosep]
  \item \textbf{Learned embedding spaces.}
    PILOT jointly optimizes query and LLM embeddings aligned by affinity,
    potentially capturing non-linear structure that our fixed
    SentenceTransformer + PCA pipeline cannot.
    GraphRouter~\cite{feng2025graphrouter} similarly learns graph-based
    representations that capture task--query--LLM interactions.
    Integrating a learned embedding layer is a natural extension.
  \item \textbf{Preference-conditioned inference.}
    BaRP allows operators to specify a performance--cost preference vector at
    inference time without retraining, providing per-request flexibility.
    banditGPT's $\lambda$ parameter achieves a similar effect but requires
    setting the value before the evaluation episode rather than per-query.
  \item \textbf{Budget-aware cost policies.}
    PILOT's knapsack formulation optimizes cumulative cost over a horizon,
    while banditGPT enforces per-request constraints.
    For deployments with strict aggregate budgets (e.g., monthly spend caps),
    a horizon-aware cost policy could complement banditGPT's per-request
    filtering.
  \item \textbf{Generalization to unseen LLMs.}
    LLM Bandit's model identity vectors and GraphRouter's inductive graph
    formulation both enable generalization to new LLMs without retraining.
    banditGPT's \texttt{register\_model()} API provides a mechanism for
    onboarding new models but requires online exploration rather than
    zero-shot generalization.
  \item \textbf{Combinatorial selection.}
    C2MAB-V~\cite{dai2024c2mab} selects \emph{subsets} of models for
    collaborative tasks, whereas banditGPT selects a single model per request.
\end{itemize}

These complementary strengths suggest fruitful directions for future
integration: a Corralling-based meta-learner with a learned embedding space
and horizon-aware cost policy would combine the safety and prior machinery
of banditGPT with the adaptive representations and budget optimization of
PILOT and BaRP.
